\section{A Responsibility Architecture for AI Deployment}
\label{sec:responsibility}

The deployment-control cascade identifies where control resides. Responsibility law must decide what to do with that information. The principal obstacle is not the absence of doctrine. Direct negligence, products liability, agency, corporate criminal liability, officer liability, discovery, preservation, contribution, and indemnity already separate organizational from individual conduct. The obstacle is sequence. An injured person is often asked to identify the internal decisionmaker before obtaining the records that reveal who decided, while the organization describes the output as an act of the system and its employees reasonably fear that candid escalation will make them the easiest defendants.

This Part proposes a two-layer architecture that changes that sequence. As an analytic baseline, the enterprise possessing material control over a deployment is its external responsibility center: once a claimant proves the elements of an applicable cause of action, the enterprise cannot defeat attribution merely because no employee composed the output or because responsibility was divided internally. For designated high-impact deployments, legislation should turn that baseline into the nondelegable duty, primary-recourse rule, and production obligation specified below. A predeployment \Charter{} supplies the internal layer: it assigns risk ownership, decision and stop authority, reporting protection, evidence duties, and recourse before the incident. The first layer preserves a practical defendant and remedy. The second distinguishes the person who warned from the person who concealed, and the person who followed a reasonable constraint from the person who bypassed it.

The front door has two components with different domains. Identification is close to declaratory where one company visibly operates an integrated service; its reforming work lies in divided deployments where several enterprises control different layers and the user cannot see which one governed the risk. Focused production matters in both settings because the deployment record remains inside the enterprise even when its name is obvious. The proposal therefore states the two obligations separately and calibrates each to the architecture that makes it necessary.

The proposal is neither blanket strict liability nor employee immunity. Substantive liability remains tied to the governing law, and personal liability remains available for a person's own actionable conduct. The architecture supplies a claimant-facing point of responsibility, a burden of producing deployment records, and a principled internal allocation.

\subsection{The Enterprise as the Claimant-Facing Responsibility Center}

An enterprise is the external responsibility center when, in the course of business or a public function, it makes the relevant AI deployment available or incorporates it into a decision process and possesses material deployment control. That means authority over one or more features capable of creating or materially reducing the risk at issue---such as purpose, eligible users, deployment data, model or vendor selection, the harness and tools, warnings, monitoring, continuation, or rollback. A direct-to-consumer provider ordinarily occupies the role itself. In an API deployment, a downstream application company, hospital, employer, or agency may be the principal deploying enterprise even though an upstream provider remains a parallel responsibility center for risks within its retained control.

Under the proposed statutory regime, designation as a ``responsibility center'' has three consequences. First, subject to the ordinary rules governing jurisdiction and immunity, the enterprise is a proper initial defendant on the statutory deployment-duty claim rather than an unnamed model or unidentified employee; in an integrated service this usually confirms an already visible party, while in a divided deployment it supplies the identification rule. Second, it bears the focused preservation and production obligation specified below. Third, internal allocation cannot by itself reduce rights that governing law gives an injured outsider. None of the three establishes defect, breach, causation, compensable harm, or a measure of damages. Those elements still perform their limiting work.

Existing doctrine already supports much of this structure, but it does not universally guarantee the proposed front door. A corporation can be directly negligent through its own design, supervision, warning, monitoring, or deployment choices. It can be vicariously liable for employee torts within the governing scope rule. A commercial seller can face product duties without locating the defect in a culpable engineer. In jurisdictions that have enacted section 14 of the Uniform Electronic Transactions Act or an equivalent rule, contracts may be formed through interactions involving authorized electronic agents. The law routinely treats an organization as the actor even when information and conduct are distributed among many employees.\lrfootnote{See \bluecite[\S\S~7.03(1), 7.05, 7.07]{restatementAgency2006}; \bluecite[\S\S~1--2]{restatementProducts1998}; \bluecite[\S~14]{ueta1999}.}

Vicarious liability alone cannot deliver the proposal because a claimant may still have to prove an employee's underlying tort and scope of employment. For designated high-impact deployments, legislatures should therefore create a nondelegable deployment duty and a primary-recourse rule for every enterprise satisfying the material-control test above with respect to the risk at issue. ``Primary recourse'' means that the enterprise is the proper defendant on the statutory deployment-duty claim and that a claimant need not identify or first proceed against an employee; it does not dispense with proof of breach, causation, legally cognizable harm, or any defense supplied by the statute. Multiple enterprises can occupy parallel lanes, with contribution and indemnity following afterward. This is the Article's normative addition, not a claim that present common law in every jurisdiction already makes the company pay first.

Agency law is illustrative. The Restatement provides that an employer is subject to liability for an employee's tort committed within the scope of employment and separately recognizes a principal's direct liability for negligent selection, training, retention, supervision, or control.\lrfootnote{See \bluecite[\S\S~7.03(1), 7.05, 7.07]{restatementAgency2006}.} Under the Third Restatement, an employee acts within scope when performing assigned work or engaging in a course of conduct subject to the employer's control; conduct leaves scope when it forms an independent course not intended to serve any employer purpose. Violation of an instruction is relevant but does not itself answer that inquiry. \emph{Ellerth} applies agency principles in the specialized Title VII supervisor-harassment setting, while \emph{Bushey} supplies a maritime enterprise-risk illustration in which unauthorized details remained characteristic of the business activity.\lrfootnote{See \casecite[756--59]{burlingtonEllerth1998}; \casecite[171--73]{bushey1968}.} AI distribution intensifies the reason for organizational attribution: the output exists only because the enterprise assembled and maintained a system intended to act without contemporaneous employee review.

Disclaimers do not change the sequence. ``AI-generated,'' ``for reference only,'' ``verify independently,'' and ``the model may make mistakes'' can be relevant warnings. Their effect depends on the claim, the clarity and timing of disclosure, assent, public policy, the audience, and whether a reasonable design could have reduced the risk. Contract law has long limited exculpation where a term conflicts with public policy or defeats a nonwaivable duty. Warranty law separately regulates disclaimer, and products doctrine asks whether a reasonable alternative design or adequate instruction would have reduced the relevant risk; a warning is not categorically a substitute for reasonable design.\lrfootnote{See, e.g., \casecite{tunkl1963}; \casecite{henningsen1960}; \statcite{ucc2316}; \bluecite[\S~2(b)--(c) \& cmt. l]{restatementProducts1998}.} No phrase creates a general license to externalize enterprise risk, just as no marketing metaphor creates an artificial defendant.

Software classification remains jurisdiction specific. United States courts have not adopted a single rule treating every stand-alone service, model, or update as a product. A useful enacted comparator is the European Union's 2024 Product Liability Directive, which expressly includes software, defines a manufacturer's control to include the ability to supply updates or upgrades, addresses substantial modification, and creates proportionate disclosure and proof rules.\lrfootnote{See \statcite{euProductLiabilityDirective2024}, arts. 4(1), 4(5), 4(18), 9--10. Member States must transpose the Directive; it neither governs a United States claim nor resolves the classification of every AI service.} Its relevance here is structural, not transplantive: retained update authority, material modification, and information access can each be assigned to an existing legal person without characterizing the software itself as one.

\subsubsection{Two Chatbots, Two Merits Results}

The Hangzhou and Air Canada disputes show how an enterprise front door can coexist with different substantive outcomes. In \emph{Moffatt}, Air Canada's website chatbot incorrectly told a customer that he could apply for a bereavement fare after travel. The airline later denied the request because its written policy required an application before travel. Before the British Columbia Civil Resolution Tribunal, the airline effectively sought to distinguish the chatbot from the rest of its website. The tribunal rejected that move: the customer was not required to determine which part of the airline's website was accurate, and Air Canada was responsible for information supplied through its channel.\lrfootnote{See \casecite[27--31]{moffatt2024}.} Attribution to the enterprise did the necessary work because the governing misrepresentation rule and evidence supported relief.

The Hangzhou Internet Court confronted more extravagant generated language. On June 29, 2025, a user asked an anonymized generative-AI application where a university program was located. After the application provided incorrect information and the user challenged it, the system apologized, purported to offer 100,000 yuan in compensation, and suggested litigation in the Hangzhou Internet Court. The user sought 9,999 yuan in damages. The court held that an AI system lacked the civil-subject status and subjective intention necessary to make the purported promise its own legal act. The text also did not manifest the provider's actual intention to pay. Treating the generative-AI offering as a service under China's Interim Measures rather than a product governed by the Product Quality Law, the court separately applied the Civil Code's general fault rule. The parties did not appeal, and the judgment was final when the court account appeared in January 2026.\lrfootnote{See \casecite{hangzhouChatbot2026} (discussing Civil Code art. 1165(1)). The publicly available court account anonymizes the application and provider and, rather than a complete published judgment, supplies the procedural and factual description used here.}

The absence of a machine promisor did not leave the wrong answer outside law. The court separately considered whether the provider had breached a differentiated, fault-based duty. The reported record included conspicuous limitation notices on the welcome page, in the user agreement, and in the interface, as well as retrieval-augmented generation to improve output reliability. It also included no adequate proof of actual loss or a causal effect on the user's application decision. On that record, the court found neither actionable fault nor compensable damage. The provider won after its conduct was examined, not because the model became the responsible party.

The two decisions therefore supply a useful pair. Air Canada used an automated channel to communicate fare information in a context where the tribunal found enterprise attribution and reasonable reliance. The anonymized Hangzhou application generated a fantastical promise outside provider manifestation, followed by a distinct negligence inquiry that the provider satisfied. A front-door rule neither binds an enterprise to every string of text nor immunizes it from every string of text. It keeps the enterprise in the case long enough for the relevant law of attribution, duty, and proof to operate.

\subsection{The Rule in Operation: \emph{Garcia}}

The narrowed rule can now be run against a complete deployment rather than stated in the abstract. \emph{Garcia} is especially revealing because the published order preserved ordinary merits questions while the pleaded stack contained an application company, two founders, Google entities, user-created characters, and a minor user's own inputs. The exercise does not decide the settled action. It shows what the front door and production rule would have required, and which disputes they would have left for ordinary law.

\subsubsection{Threshold and Initial Responsibility Center}

The threshold would have been satisfied on the allegations accepted for Rule 12 purposes. The complaint pleaded recognized injuries, including wrongful death, and a material nexus to months of interaction with the Character.AI deployment. It targeted design and distribution features---including anthropomorphic presentation, access by minors, relational interaction, warnings, and asserted safety omissions---rather than naming the generated character as a defendant. The district court held the application plausibly alleged to be a product and allowed core product and negligence theories to proceed beyond pleading.\lrfootnote{See \casecite[1165--69, 1179--83]{garcia2025}. The court did not find that the deployment caused Sewell's death or that any challenged feature was defective.}

Character Technologies would be the clearest initial responsibility center. It made the consumer deployment available and, on the pleaded record, operated the application through which users selected or created characters and interacted with generated responses. The risk asserted was tied to features an application operator ordinarily can govern: eligibility, persona presentation, memory and continuity, interaction policies, crisis detection, warnings, monitoring, and continued access. This is the vertically integrated setting in which the identification component adds little. The value of the front door lies in confirming that an \term{automatic instrumentality} does not end attribution merely because no employee selected the final exchange, and in opening the production obligation without requiring the claimant to identify an engineer.

The other named actors illustrate why ``responsibility center'' must remain risk specific. The complaint alleged that the founders personally designed, directed, and released the technology and that Google supplied substantial technical, financial, and organizational assistance with relevant knowledge.\lrfootnote{See First Amended Complaint paras. 24--25, 50--72, \emph{Garcia v. Character Technologies, Inc.}, No. 6:24-cv-01903-ACC-EJK, Dkt. 11 (M.D. Fla. Nov. 9, 2024) [hereinafter \emph{Garcia FAC}]; \casecite[1165--66, 1183--84]{garcia2025}.} Those allegations may support individual direct or aiding-and-abetting theories under governing law. They would not make an investor, former employer, licensor, or infrastructure supplier a deployment responsibility center by title alone. The statutory inquiry would ask whether each enterprise retained material authority over the relational risk---for example, control of the relevant model, application configuration, safety information, release condition, update, or stopping decision. The production record would test that proposition rather than presume it from corporate association.

\subsubsection{The Focused Production Record}

After the threshold, Character Technologies would identify the deployed model and application versions, the entities occupying the cascade, and the nonprivileged portions of the charter and incident record proportionate to the pleaded risk. In this case, the focused record would include the age-access rule and its enforcement; the character definition and features attributable to its creator; system instructions governing relational and self-harm interactions; memory and proactive-engagement settings; warnings actually displayed to this user and their timing; relevant evaluation protocols and results; crisis detections and interventions; material complaints and escalations received before the events; release approvals; and available restriction, rollback, or account-intervention controls.

That is not a demand for every model weight, every user's conversation, or unrestricted disclosure of a minor's data. Rule 26 proportionality, privilege, privacy protections, trade-secret orders, redaction, and staged discovery would continue to govern. A first stage could identify the system, responsible entities, governing policies, and preserved incident-specific events. Broader evaluation or design records would follow only as the pleaded theory and initial production justified them. The rule converts ``the algorithm'' into bounded requests tied to the control variables in Part I.

The record would also preserve facts favorable to the defense. The pleadings and motion papers disputed the effect of user selection, role play, user-edited messages, platform rules, content blocking, and an in-chat notice that character statements were fictional.\lrfootnote{See \casecite[1167--69]{garcia2025}.} A configuration and event record could show which character features came from a third-party creator, which messages the user edited, which interventions fired, what warnings appeared, and whether a proposed safeguard would have altered the sequence. Production replaces categorical claims of unpredictability with evidence capable of supporting either side.

\subsubsection{From Record to Merits and Internal Allocation}

With that record available, ordinary doctrine would do the remaining work. Product classification would determine which design and warning rules applied. The process model would organize what was foreseeable at the relevant time; feedback would identify notice; authority and resources would locate feasible precautions; and the event sequence would test cause in fact and legal causation. The empirical studies discussed in Part III identify variables that providers can now test, but studies published after the 2023--2024 interactions cannot by themselves establish what a defendant historically knew or what precaution was then feasible. The proper inquiry is temporal and deployment specific.

The same discipline applies to later safety measures. A post-incident self-harm intervention may demonstrate that a control is technically describable, while evidence law and the substantive standard determine whether it is admissible or proves a pre-incident alternative.\lrfootnote{See Fed. R. Evid. 407; \casecite[1169]{garcia2025}.} The architecture neither converts subsequent remediation into liability nor permits the absence of contemporaneous records to be concealed behind a general claim that model behavior cannot be reconstructed.

Finally, the charter would separate external responsibility from individual recourse. It would show who approved minor access, who owned relational and self-harm risks, what dissent reached management, who controlled release resources, and whether anyone concealed an incident or bypassed a constraint. A documented warning by a safety employee would support protection, not an inference that the employee should become the claimant's sole target. A knowing suppression or unauthorized removal would remain available for individual consequences. Because no statutory charter governed the historical deployment, this application is prospective: it shows the allocation and evidence the proposed rule would require for the next high-impact relational system.

\emph{Garcia} thus demonstrates the architecture's exact payoff. The initial enterprise was visible, so identification was easy. The hard questions concerned the record behind the interface, the division of retained control, and the causal significance of particular design choices. The production rule reaches those questions without predetermining their answers. The settled case remains a classification and pleading anchor; the worked example supplies the procedure that could have carried its merits analysis forward.

\subsection{Corporate and Individual Responsibility Run in Parallel}

Three doctrines are easily conflated. \term{direct enterprise responsibility} concerns the organization's own design, policy, resource, supervision, warning, and deployment choices. \term{vicarious responsibility} attributes a person's conduct to an organization under agency or another relationship. \term{individual direct responsibility} concerns a person's own tortious, criminal, or otherwise actionable conduct. The three are analytically distinct. Direct enterprise and individual liability can arise independently; vicarious liability ordinarily presupposes an underlying employee tort even when the employee need not be named or personally recoverable.

The separation matters in both directions. A corporation does not avoid its own negligence by identifying an employee who implemented the policy. An employee does not become liable merely because the organization is liable. But acting for a corporation is not a general immunity from one's own tort. Restatement (Third) of Agency section 7.01 states that an agent is subject to liability to a third party harmed by the agent's tortious conduct, whether or not the principal is also liable.\lrfootnote{\bluecite{restatementAgency2006}.} The claimant must still prove that individual's duty, conduct, causation, and any required state of mind.

Veil piercing answers another question. It permits exceptional disregard of entity separateness where owners misuse the corporate form under jurisdiction-specific standards. It is not ordinarily required to sue an officer, engineer, clinician, or manager for that person's own tort.\lrfootnote{See \casecite[503--05]{francesT1986}; \casecite[1380--81]{pmcKadisha2000}. Position or general supervisory power alone remains insufficient; the claimant must establish personal participation, authorization, duty, and causation under the applicable rule.} Treating the two as synonyms either overprotects individual wrongdoing or threatens routine employees with owners' liabilities. The charter should maintain the line: to the extent applicable law permits, ordinary good-faith work is allocated and indemnified internally; concealment, falsification, and personally tortious action are evaluated directly.

Corporate criminal law confirms that organization and individual can answer in parallel. \emph{New York Central \& Hudson River Railroad Co.\ v. United States} upheld attribution of employee conduct to a corporation for a federal offense committed within authority and for corporate benefit.\lrfootnote{\casecite{newYorkCentral1909}.} The responsible-corporate-officer doctrine, in tightly regulated public-welfare settings, can impose liability on an officer with responsibility and authority to prevent or correct violations even without conventional participation in the prohibited act.\lrfootnote{See \casecite{unitedStatesPark1975}.} State positive law can also authorize corporate liability for negligent homicide. In \emph{Commonwealth v. Angelo Todesca Corp.}, the Massachusetts Supreme Judicial Court upheld a corporation's conviction under that State's motor-vehicle-homicide statute for its driver's negligent operation while conducting authorized corporate business.\lrfootnote{See \casecite[133--37]{todesca2006}.} That jurisdiction-specific rule should no more be generalized into a national AI offense than \emph{Park} should be generalized into automatic executive liability. Together, the authorities demonstrate a structural possibility: organizational liability need not displace individual responsibility, and managerial authority matters where positive law makes it matter.

Corporate fiduciary doctrine supplies a further, carefully bounded analogy. \emph{Caremark} and its successors ask whether directors made a good-faith effort to establish board-level information and reporting systems for mission-critical compliance risk.\lrfootnote{See \casecite{caremark1996}; \casecite{marchand2019}; \casecite{boeingDerivative2021}.} These are fiduciary claims belonging to the corporation, not a free-standing tort duty owed by directors to every person affected by AI. Their relevance is institutional. Law already recognizes that an organization can fail through the absence of reporting architecture, escalation, and board attention even when no director designed the failed component.

The charter converts that insight into a deployment-specific record. It allows a court to distinguish at least four situations: management declined resources after receiving a documented risk; a deployer acted within an approved and reasonably monitored design; an employee raised a concern that was suppressed; or a person knowingly bypassed a constraint. Without an ex ante allocation, all four can look like undifferentiated ``AI failure.''

\subsection{The Mandatory Responsibility Charter}

For designated high-impact deployments, law should require a written \Charter{} before release. ``High impact'' should be defined by the governing sector or enactment---for example, systems that materially determine access to employment, credit, housing, education, health care, liberty, essential public benefits, or physical safety, and relational systems offered to minors. A general-purpose model does not become high impact merely because it could be used in those settings; the obligation attaches to the covered deployment and to providers only for functions the sectoral rule assigns them.

The charter is not an aspirational ethics statement. It is a control instrument tied to a named deployment, version, and operating context. At minimum, it should contain the following:

\begin{longtable}{>{\raggedright\arraybackslash}p{0.25\textwidth} >{\raggedright\arraybackslash}p{0.33\textwidth} >{\raggedright\arraybackslash}p{0.32\textwidth}}
\caption{Minimum Contents of a Responsibility Charter}\label{tab:charter}\\
\toprule
Component & Required allocation & Evidentiary function \\
\midrule
\endfirsthead
\toprule
Component & Required allocation & Evidentiary function \\
\midrule
\endhead
Deployment identity and scope & System version, purpose, population, data, tools, vendors, material limits, and change threshold. & Defines the system at issue and distinguishes later material modification. \\
Named risk owners & A responsible executive and an operational owner for each material hazard, with adequate authority and resources. & Identifies who receives information and can commission mitigation. \\
Deployment sign-off & Specified approvers, evidence considered, unresolved dissent, conditions, and expiration or review date. & Records authorization, knowledge, alternatives, and limits. \\
Protected escalation and dissent & Channels outside the ordinary reporting line; direct access to an independent reviewer or governing body; response deadlines. & Preserves notice, prevents suppression, and separates warning from concealment. \\
Stop and rollback authority & Who may pause, restrict, revoke credentials, restore a version, or decommission; triggers and response times; tested procedures. & Establishes practical ability to prevent or mitigate and tests whether oversight is real. \\
Incident and configuration records & Versions, prompts and policies where proportionate, tool calls, alerts, overrides, complaints, decisions, and corrective actions; retention schedule. & Makes attribution, causation, discovery, and learning possible while allowing privacy minimization. \\
Independent safety review & Reviewer independence, access to data, evaluation scope, conflict rules, and remediation tracking. & Tests self-assessment and documents feasible controls. \\
Internal allocation and recourse & Indemnification, discipline, contribution, and safe-harbor standards fixed ex ante. & Protects good-faith compliance and preserves recourse for knowing misconduct. \\
External-rights savings clause & Express statement that the charter does not waive or narrow claimant rights, statutory duties, or regulatory powers. & Prevents internal governance from becoming an exculpatory contract with outsiders. \\
\bottomrule
\end{longtable}

The components have distinct institutional pedigrees. Named owners, adequate authority and resources, and documented approval track the Organizational Sentencing Guidelines, FAA accountable-executive rules, and NIST's assignment of lifecycle and executive responsibility.\lrfootnote{See U.S. Sent'g Guidelines Manual \S~8B2.1(b)(2), \bluecite{ussgCompliance2025}; \statcite{faaSMSPart5}, \S~5.25; \bluecite{nistAIRMF2023}, GOVERN 2.1, 2.3.} Protected escalation tracks compliance-program reporting channels and sectoral anti-retaliation rules.\lrfootnote{See U.S. Sent'g Guidelines Manual \S~8B2.1(b)(5)(C), \bluecite{ussgCompliance2025}; \statcite{energyReorgWhistleblower}.} Stop and rollback authority and configuration records map respectively to NIST's decommissioning and post-deployment response functions and the EU AI Act's intervention and logging requirements; continuing assessment, assurance, and recordkeeping map to FAA safety-management rules. Independent-review safeguards draw instead on ex ante audit and disclosure proposals designed to preserve investigation incentives.\lrfootnote{See \bluecite{nistAIRMF2023}, GOVERN 1.7, MANAGE 4.1; \statcite{euAIAct2024}, arts. 12, 14(4)(e), 17--20; \statcite{faaSMSPart5}, \S\S~5.53, 5.71--5.73, 5.97; \artcite[1047--50]{ramakrishnanTortFrontier2026}.} Internal recourse has a separate genealogy in conditional indemnification and compliance enforcement, while the external-rights savings clause preserves the public-policy limits on private exculpation.\lrfootnote{See \statcite{dgcl145}; U.S. Sent'g Guidelines Manual \S~8B2.1(b)(6), \bluecite{ussgCompliance2025}; \casecite{tunkl1963}.}

This is not a governance scheme invented from whole cloth. The Organizational Sentencing Guidelines define an effective compliance program through leadership responsibility, operational responsibility with adequate resources and authority, training, monitoring, reporting without fear of retaliation, enforcement, response, and periodic risk assessment.\lrfootnote{See U.S. Sent'g Guidelines Manual \S~8B2.1, \bluecite{ussgCompliance2025}.} The Department of Justice's compliance-program guidance asks whether a program is well designed, adequately resourced and empowered, and effective in practice; its 2024 update specifically asks how a company governs AI risks, controls unintended consequences, and addresses deliberate or reckless misuse.\lrfootnote{See \bluecite{dojECCP2024}.} These materials arise in criminal enforcement and sentencing. They are transferable here not as liability rules, but because they operationalize information, authority, incentives, and response inside complex enterprises.

Aviation supplies a stronger safety-system analogy. On a staged implementation schedule, federal rules require specified aviation organizations to develop and maintain an organization-wide safety management system with accountable executive responsibility, hazard identification, risk controls, assurance, training, communication, and records.\lrfootnote{See \statcite{faaSMSPart5}, \S\S~5.1--5.17, 5.21--5.97; \bluecite{faaSMS2026}.} The Aviation Safety Reporting Program seeks candid safety information through third-party reporting and a qualified, nonpunitive enforcement policy rather than blanket immunity.\lrfootnote{See \statcite{aviationReporting}.} The analogy is functional: AI deployments, like aviation systems, couple technical components with operators and organizational decisions, and near-miss information is easily suppressed if disclosure brings only personal exposure.

AI and medical governance sources converge on the same elements. NIST calls for documented lifecycle roles, executive responsibility, external feedback, risk-prioritized resources, third-party controls, monitoring, and safe decommissioning.\lrfootnote{See \bluecite{nistAIRMF2023}.} A 2026 consensus produced by more than forty Chinese medical and research institutions organizes medical-AI governance around access evaluation, clinical application, patient rights, data governance, risk management, and competence; it specifically identifies multidisciplinary review, real-world validation, traceability, dynamic monitoring, human--machine responsibility, and circuit breakers.\lrfootnote{See \artcite{caoMedicalConsensus2026}.} Recent medical scholarship similarly identifies four conditions for meaningful oversight---epistemic capacity, cognitive space, decisional authority, and intervention effectiveness---and treats oversight as a property of the clinical system rather than a clinician's ceremonial presence.\lrfootnote{See \artcite{vanDeSandeOversight2026}.}

The convergence is the point. Criminal compliance, aviation safety, NIST risk management, and medical governance serve different ends and should not be collapsed into one legal standard. They nonetheless operationalize overlapping elements: named responsibility, information flow, authority, resources, protected reporting, review, and corrective capacity. Recent tort scholarship reaches the same institutional need from the opposite direction: liability by itself can discourage risk investigation and disclosure, while ex ante rules, independent auditors, and disclosure obligations can preserve the information on which prevention and adjudication depend.\lrfootnote{See \artcite[1047--50]{ramakrishnanTortFrontier2026}.} The charter packages that shared infrastructure into a document a deployment and a court can use.

Two limits prevent compliance theater. First, possession of a charter is not compliance. Courts and regulators must ask whether the named people had actual authority and resources, reports reached them, stop mechanisms worked, and material changes triggered review. Second, unless governing positive law provides otherwise, compliance should be evidence rather than a conclusive defense. A carefully completed charter can support reasonableness; it cannot nullify a statutory duty, excuse a known unreasonable risk, or bar proof that the system operated outside its stated scope.\lrfootnote{See \bluecite[\S~16(a)]{restatementPhysicalEmotional2010}; \bluecite[\S~4(b)]{restatementProducts1998}. Jurisdictions differ, and a legislature or valid preemption rule may assign compliance a different effect.}

\subsection{Who Can Enact and Administer the Charter}

The proposal has no mystery legislature. Its designated-deployment structure permits adoption by the lawmaker that already governs the underlying activity. State legislatures can attach the charter, nondelegable duty, primary recourse, and production rule to companion services, health care, employment, education, insurance, consumer transactions, or other fields within state authority. Congress can legislate for federal programs, procurement, interstate sectors, and national schemes within its enumerated powers. Either lawmaker can authorize a sectoral regulator to specify forms, reporting intervals, change thresholds, and technical controls within intelligible statutory boundaries. The enacting law should identify the covered deployment, duty holder, minimum charter contents, production trigger, remedy, and relation to overlapping law; administration can then remain technically adaptive without relocating those policy choices.

California's companion-chatbot law demonstrates the form at state scale. Senate Bill 243 defines the covered relational service and its operator, requires conspicuous nonhuman disclosure in specified circumstances, mandates a published self-harm protocol and protections for known minors, requires reporting to the Office of Suicide Prevention beginning in 2027, and supplies a private remedy while preserving cumulative law.\lrfootnote{See \statcite{californiaSB2432025}.} It does not enact this Article's charter: it does not require named risk owners, deployment sign-off, protected escalation, stop authority, or the litigation production rule. It proves the adjacent institutional point. A legislature can designate a relational deployment, name the operator, combine protocol and reporting duties, and attach enforceable consequences without defining every general-purpose model as high impact.

Federal practice supplies implementation tools with different legal effects. In its 2025 section 6(b) companion-chatbot study, the Federal Trade Commission ordered named firms to report organizational roles, predeployment standards and thresholds, the positions of deployment approvers, considered and rejected mitigations, post-deployment monitoring, complaints, employee concerns, escalation, and response.\lrfootnote{See \bluecite{ftcCompanionInquiry2025}; \bluecite{ftcCompanionOrder2025}, specifications 3, 11--17.} The compulsory study is neither a violation finding nor a generally applicable substantive rule. It shows that the organizational and deployment record the charter would standardize can be specified, collected, and verified rather than reconstructed as an abstract demand for ``the algorithm.''

The Department of Justice's 2024 compliance guidance works through enforcement evaluation rather than rulemaking. It asks prosecutors whether a company has assessed emerging-technology risk, governed business and compliance uses of AI, controlled unintended consequences and deliberate misuse, established accountability and human baselines, enabled confidential reporting, and continuously monitored and corrected model behavior.\lrfootnote{See \bluecite[3--4, 7, 18]{dojECCP2024}.} The guidance does not create an ex ante tort duty. Its significance is operational: federal corporate-enforcement practice already evaluates the authority, information flow, incentives, and correction mechanisms the charter records.

FAA Part 5 supplies the mature delegated model. Its staged schedule requires covered aviation organizations to develop and, when applicable, maintain a formal, organization-wide safety management system with an accountable executive, identified management authority to accept safety risk, hazard analysis, corrective action, and retained records. For organizations subject to section 5.71(a)(7), the system also includes confidential employee reporting without fear of reprisal.\lrfootnote{See \statcite{faaSMSPart5}, \S\S~5.3, 5.9(e), 5.11--5.17, 5.21, 5.23, 5.25, 5.53--5.57, 5.71--5.75, 5.95--5.97.} Part 5 is not an AI liability rule and does not supply every charter component. It demonstrates that named authority, protected reporting, continuing assurance, and evidence retention can be administered as an operating system rather than an aspirational policy.

This division of labor answers the concern created by \emph{Loper Bright}. An agency may implement a charter only within authority Congress or a state legislature has conferred; technical expertise does not itself create the primary-recourse duty or private remedy.\lrfootnote{See \casecite[394--96]{loperBright2024}.} Courts, meanwhile, can use present law now: they can distinguish direct from vicarious responsibility, apply ordinary attribution, enforce preservation, order proportional discovery from proper parties and nonparties, and refuse to treat an AI system as an assetless substitute for a defendant. They cannot create the proposed statutory front door by discovery order alone. Legislative adoption supplies the duty and trigger; bounded agency implementation supplies sector-specific administration; adjudication supplies independent interpretation and enforcement.

\subsection{Information Asymmetry and the Production Rule}

The claimant-facing layer fails if the enterprise can identify itself yet withhold the control record. AI deployments produce unusually concentrated information: model and application versions, system instructions, retrieval sources, tool permissions, evaluation results, incident queues, escalation tickets, internal approvals, and rollback events. The injured person usually sees only an output or decision. The enterprise knows which of several vendors, models, and configurations was active.

Ordinary civil procedure provides the foundation. Rule 26 permits discovery of relevant, proportional nonprivileged matter and requires initial disclosure of information and documents a party may use to support claims or defenses.\lrfootnote{\statcite{frcp26}.} Once litigation is reasonably anticipated, preservation duties attach under governing law. Rule 37(e), which draws on but does not itself create that duty, addresses electronically stored information that should have been preserved, is lost because reasonable steps were not taken, and cannot be restored or replaced. Curative measures require prejudice; an adverse-inference instruction or equivalent severe sanction requires intent to deprive.\lrfootnote{\statcite{frcp37e} advisory committee's note to 2015 amendment.} The charter should make reasonable preservation possible before a court has to reconstruct an undocumented system.

For covered high-impact deployments, legislation should add a focused \term{production rule}. Once a claimant pleads a recognized injury and a plausible material nexus to the deployment, the responsibility center must identify the system and entities in the cascade and produce the nonprivileged portions of the charter and incident record proportionate to the claim. Sensitive model, security, health, and personal information can be redacted or protected by order, including an order limiting disclosure of trade secrets or confidential research and commercial information.\lrfootnote{See \statcite{frcp26}, r. 26(c)(1)(G). The EU Product Liability Directive offers a comparative enacted model of proportionate disclosure, protection of confidential information, and bounded presumptions; it is not current United States procedure. See \statcite{euProductLiabilityDirective2024}, arts. 9--10.} Secrecy is not a reason to leave the system unidentified. The production obligation is procedural. It does not shift the ultimate burden of proving breach or causation unless the governing legislature expressly says so.

Existing evidentiary doctrine demonstrates why a calibrated inference may sometimes be warranted. Section 3 of the Restatement (Third) of Torts: Products Liability permits an inference that a defect existed at the time of sale or distribution, without proof of a specific defect, when the incident is of a kind that ordinarily results from product defect and was not solely the result of causes other than such a defect.\lrfootnote{\bluecite[\S~3]{restatementProducts1998}; see also \casecite[41--43]{speller2003} (applying New York law and requiring evidence that excludes other causes not attributable to defendants).} Res ipsa loquitur likewise permits a factfinder in defined circumstances to infer negligence from the type of event. Jurisdictions formulate the doctrine differently; modern formulations do not invariably treat rigid exclusive control as a separate element.\lrfootnote{See \bluecite[\S~17 \& cmt. b]{restatementPhysicalEmotional2010}.} These doctrines are bounded; neither says that any unexpected AI output proves negligence. The charter instead makes specific proof more available and reserves inference for the jurisdiction's ordinary conditions.

A missing record should have equally bounded consequences. If the system never captured information that no duty required it to capture, the absence is not automatically spoliation. In federal court, Rule 37(e) governs measures for ESI lost after a litigation-preservation duty arose and forecloses reliance on inherent authority or state law to prescribe different measures for that loss; it does not displace an otherwise applicable independent state-law spoliation claim.\lrfootnote{See \statcite{frcp37e} advisory committee's note to 2015 amendment.} If a sectoral charter mandate required logs and the enterprise omitted them, the statute should specify the consequence; violation can be evidence without becoming automatic liability for the underlying harm. This calibration makes documentation valuable rather than punitive by default.

\subsection{Protect the Reporter; Pursue the Concealer}

An accountability system that directs every loss toward the nearest engineer will suppress the information it needs. The charter should therefore establish an internal safe harbor for a worker who, in good faith and on an objectively reasonable basis, documents a material risk, uses an authorized escalation or stop channel, preserves evidence, and acts within an approved role. The default internal consequences should be defense and indemnification by the enterprise to the extent permitted by governing law, protection against retaliation, and no adverse employment action based solely on protected reporting or justified intervention. Corporate indemnification statutes supply a model but also impose conditions; the charter cannot promise indemnity that positive law forbids.\lrfootnote{See \statcite{dgcl145}.}

The proposed protection for refusal should be comparably defined. A worker receives it when the worker reasonably believes the requested act is unlawful or presents an imminent risk of death or serious physical harm, gives the enterprise notice where feasible, and lacks time or a reasonably effective ordinary channel before the risk materializes. This is a proposed threshold informed by, not identical to, sectoral law.\lrfootnote{Compare \statcite{energyReorgWhistleblower} (protecting specified reports and refusal to engage in an identified unlawful practice), with \casecite[9--13]{whirlpool1980} (upholding a narrow occupational-safety protection for good-faith refusal in the face of reasonably perceived imminent danger and inadequate time for ordinary redress).} A statute should also specify causation and remedy; the Energy Reorganization Act's contributing-factor test and clear-and-convincing same-action defense provide one balanced model.\lrfootnote{See \statcite{energyReorgWhistleblower}, \S~5851(b)(3).}

The anti-retaliation component has a substantial statutory genealogy. The Sarbanes--Oxley Act protects specified reports of fraud and securities violations; the False Claims Act protects efforts to stop violations; and the Occupational Safety and Health Act protects safety complaints and related participation.\lrfootnote{See \statcite{soxWhistleblower}; \statcite{falseClaimsRetaliation}; \statcite{oshaRetaliation}.} Aviation reporting programs encourage voluntary disclosure of safety events. The Patient Safety and Quality Improvement Act separately protects good-faith reporters, but its broad privilege and confidentiality rules for patient-safety work product should not be imported into this proposal because they would conflict with the focused production rule.\lrfootnote{See \statcite{aviationReporting}; \statcite{psqia2005}, \S~299b-22(a)--(b), (e).} The statutes' scopes, procedures, and privileges differ. They establish the transferable principle that accurate internal information is a safety resource and that retaliation can destroy it, not a general immunity from otherwise applicable law.

The safe harbor should have equally clear exceptions. It does not protect intentional or reckless infliction of harm, knowing falsification or destruction of records, concealment of a material incident, retaliation against another reporter, corruption of an evaluation, or unauthorized removal of a safety constraint. Nor does internal defense or indemnification eliminate liability that law independently imposes on the person's conduct. The charter should identify the standard and decisionmaker before a conflict arises, with independent review when senior management is implicated.

This allocation corrects an incentive mismatch. A safety engineer often has the clearest view of a failure mode but not the budget or release authority to fix it. A manager may control the schedule and resources but not understand the model detail. If both face an undifferentiated threat, the engineer may remain silent and the manager may avoid written decisions. Protecting documented escalation while preserving recourse against suppression makes authority and accountability coincide.

\subsection{The Two-Layer Rule}

The proposed statutory rule can be stated in five steps.

\begin{enumerate}[leftmargin=2.2em]
\item A claimant pleads a recognized injury and a plausible material nexus to a covered high-impact AI deployment.
\item Each enterprise that makes the covered deployment available or incorporates it into a decision process and possesses material deployment control over the risk at issue is an external responsibility center. Multiple enterprises may qualify. None can defeat the statutory claim solely by showing that no natural person selected the precise output.
\item Each responsibility center identifies the deployed system and the entities occupying the relevant cascade, then produces the proportionate, nonprivileged charter and incident record. Identification confirms an already visible operator in an integrated deployment and exposes divided control where the stack is hidden; production supplies the merits record in both. Discovery, privilege, protection orders, and preservation law govern implementation.
\item With that record available, the claimant retains the ultimate burden to prove the elements of the governing cause of action unless positive law provides otherwise. The charter neither waives outsider rights nor conclusively establishes reasonable care; it governs indemnification, defense, discipline, contribution, and recourse, and supplies evidence of authority, knowledge, and conduct.
\item Good-faith reporting, compliance, and justified intervention receive the protection specified by the statute and charter; knowing concealment, falsification, retaliation, and unauthorized bypass remain available grounds for internal discipline and for consequences independently authorized by positive law.
\end{enumerate}

The rule resolves four problems at once. A victim has an identified enterprise against which a recognized claim can proceed if its elements are proved, without first reverse-engineering an org chart. A conscientious engineer can create a safety record without volunteering as the sole loss bearer. Management responsibility follows budget, release, and response authority rather than rhetorical commitment. And organizational fragmentation becomes evidence in a mapped cascade rather than a defense based on no single person's complete control.

Legislatures should mandate the charter and production rule in defined high-impact sectors. Courts can apply the architecture's underlying distinctions now: separate direct from vicarious liability; reject a machine as a liability sink; use ordinary attribution rules; enforce preservation duties; and evaluate individual conduct without demanding veil piercing. Once jurisdiction and the ordinary party or nonparty predicates are satisfied, courts can order proportional discovery from an entity possessing, controlling, or holding the relevant records; they cannot create the proposed statutory front door by discovery order alone.\lrfootnote{See \statcite{frcp26}; \statcite{frcp34and45}.} Regulatory adoption would make the allocation uniform and ex ante.

\subsection{Residual Responsibility and Rebuttal}

An enterprise-facing rule must leave room for a genuinely different control path. Six recurring defenses illustrate how the cascade disciplines rather than erases rebuttal.

\paragraph{Intentional user circumvention.}
A user who deliberately jailbreaks a system, supplies false context, defeats a reasonably designed confirmation, or directs a prohibited act may bear comparative or primary responsibility under governing law. The enterprise can show the constraint, warning, foreseeable bypass testing, and causal effect. The word ``jailbreak'' alone is not a defense if the technique was common, encouraged by the product's use, or readily preventable.

\paragraph{Modification.}
A downstream actor who removes safeguards, changes weights, adds dangerous tools, or deploys outside documented limits may become the relevant controller. New York law illustrates the distinction: substantial modification can defeat a design-defect theory while a failure-to-warn claim based on foreseeable modification may remain available in appropriate circumstances.\lrfootnote{See \casecite[237--43]{liriano1998}.} Other jurisdictions define modification and warning duties under their own statutes and common law. Version and provenance records should show the change rather than leaving the issue to brand attribution.

\paragraph{Cyberattack.}
Unauthorized intrusion may be an intervening cause, but security is itself a foreseeable deployment condition. The questions are which threat was reasonably in scope, whether credentials and tools were exposed, which alerts arrived, and whether response was possible. A sophisticated novel attack and an unpatched known vulnerability do not occupy the same control position.

\paragraph{Government compulsion.}
A contractor may invoke the federal military-equipment defense only where the United States approved reasonably precise specifications, the equipment conformed to them, and the supplier warned the United States of dangers known to the supplier but not the government.\lrfootnote{See \casecite[512]{boyle1988}.} Two 2026 decisions sharpen the control inquiry. \emph{GEO Group v. Menocal} holds that \emph{Yearsley} supplies a merits defense, not derivative sovereign immunity, for conduct the government lawfully authorized and directed.\lrfootnote{\casecite[444--48]{geoGroupMenocal2026}.} \emph{Hencely v. Fluor Corp.} rejects broader displacement where the government neither ordered nor authorized the challenged conduct: \emph{Boyle} requires a significant conflict with an identifiable federal policy and protects a contractor only insofar as the suit challenges a government decision the contractor carried out.\lrfootnote{\casecite[37--42, 46]{hencely2026}.} A procurement relationship, political pressure, silence about the challenged means, or conduct contrary to instructions does not transfer control by itself. Part VI addresses the reciprocal limit on governmental coercion.

\paragraph{Independent deployer override.}
An upstream provider may rebut responsibility for a harm tied to a downstream purpose, data set, tool, or override it neither controlled nor reasonably could prevent. Conversely, contractual reservation of a right to audit, update, suspend access, or dictate safeguards is evidence of retained control. The result should follow the deployed function, not a standardized label that declares every party an independent contractor.

\paragraph{Unavoidable event.}
Some harmful outputs will remain unforeseeable after reasonable testing, monitoring, warning, and response. Ordinary standards of fault leave room for that conclusion; the published account of the Hangzhou decision supplies a comparative illustration rather than controlling United States law. A charter strengthens the defense by showing what risk was assessed, why controls were selected, what feedback existed, and why the proposed alternative would not have prevented the injury.

Residual responsibility therefore does not undermine the architecture. It is what makes the architecture a legal framework rather than a slogan. The enterprise faces the claimant and produces the map; the parties then litigate where the causal control path actually ran under the governing burdens of proof. The machine never has to be invented as the missing defendant.
